\documentclass[11pt,a4paper]{article}
\usepackage{DDTA_DERMATRIAGE_GUIDED_REVIEW_STYLE_R1}
\providecommand{\tightlist}{\setlength{\itemsep}{0pt}\setlength{\parskip}{0pt}}
\begin{document}
\selectlanguage{italian}
\begin{titlepage}
\thispagestyle{empty}
\vspace*{10mm}
{\sffamily\bfseries\color{DDTANavy}\fontsize{15}{18}\selectfont Documentation-Driven Threat Analysis}\par
\vspace{5mm}
{\sffamily\bfseries\color{DDTANavy}\fontsize{27}{31}\selectfont DermaTriage Extended Documentation Case Study}\par
\vspace{2mm}
{\sffamily\color{DDTABlue}\fontsize{17}{21}\selectfont Annotated Holdout Study Through A5 Closure}\par
\vspace{10mm}
\begin{tcolorbox}[enhanced,arc=3pt,boxrule=0pt,colback=DDTANavy,left=12pt,right=12pt,top=10pt,bottom=10pt]
{\color{white}\sffamily\bfseries Current research artifact}\par
{\color{white!88}\small Repository baseline: \texttt{0e60754d21aa24ea487f3f60803b6b0cce8d2e2b}. Methodology authority remains R4 unless this document explicitly states otherwise.}
\end{tcolorbox}
\vfill
{\sffamily\small\color{DDTAGray} 5 settembre 2026}
\end{titlepage}
\tableofcontents
\clearpage
\section{DermaTriage - Extended DDTA Documentation Case
Study}\label{dermatriage---extended-ddta-documentation-case-study}

\textbf{Status:} ANNOTATED HOLDOUT STUDY / THROUGH A5 CLOSURE\\
\textbf{Purpose:} preserve the reasoning, rejected interpretations, gaps
and methodology observations that must not contaminate the clean project
documentation\\
\textbf{Repository baseline:}
\texttt{0e60754d21aa24ea487f3f60803b6b0cce8d2e2b}\\
\textbf{Methodology authority:}
\texttt{DDTA\_DOCUMENTATION\_BA\_AUTHORING\_GUIDE\_R4}\\
\textbf{Source package SHA-256:}
\texttt{E9ED2C507BEFB95F54A52084687CD1E8798863AE81CF69D09568864D8CBF280E}

\subsection{1. Why this case study
exists}\label{1-why-this-case-study-exists}

DermaTriage is used as a structurally different holdout after the Facial
Access case. The goal is not to make DermaTriage look complete. The goal
is to test whether DDTA can reconstruct governed project meaning from
ordinary architecture/test/configuration documentation while:

\begin{itemize}
\tightlist
\item
  separating responsibility from technology;
\item
  separating project commitments from current implementation;
\item
  preserving source gaps;
\item
  avoiding BA contamination;
\item
  producing human-readable documentation that can later support security
  analysis.
\end{itemize}

The clean project candidate and this study are intentionally separate.
The clean candidate answers \textbf{what the project currently governs}.
This study answers \textbf{how and why that meaning was derived}.

\subsection{2. Source set and evidence
roles}\label{2-source-set-and-evidence-roles}

Authoritative package:

\texttt{DermaTriage-Docs-20260830T152637Z-1-001.zip}

SHA-256:

\texttt{E9ED2C507BEFB95F54A52084687CD1E8798863AE81CF69D09568864D8CBF280E}

Contained artifacts:

\begin{itemize}
\tightlist
\item
  \texttt{OR2\_Architecture\_Document.pdf} - principal semantic source;
\item
  \texttt{OR2\_Model\_Test\_Report.pdf} - model/evaluation evidence;
\item
  \texttt{OR3\_Dataset\_Metadata\_Catalog.pdf} -
  dataset/privacy/current-data evidence;
\item
  \texttt{OR4\_Training\_Environment\_Config.pdf} -
  configuration/realization evidence;
\item
  \texttt{OR4\_Training\_Cycles\_Report.pdf} -
  retraining/evaluation/lifecycle evidence;
\item
  \texttt{OR5\_Test\_Environment\_Setup.pdf} - test/API/environment
  evidence;
\item
  \texttt{efficientnet\_b4.pth} - model artifact, realization evidence.
\end{itemize}

Old DDTA reconstructions and old BA artifacts were excluded from
project-authority decisions.

\subsection{3. A0 - Authority gate}\label{3-a0---authority-gate}

\subsubsection{Question}\label{question}

Which documents may define DermaTriage project meaning?

\subsubsection{Result}\label{result}

The original package is authority. OR2 Architecture is the principal
semantic source because it describes project purpose, pipeline,
adaptation layer, external service interactions and self-learning loops.
Tests/training/configuration documents can corroborate or expose current
realization but do not automatically convert implementation details into
normative requirements.

\subsubsection{Methodology lesson}\label{methodology-lesson}

A traditional architecture document often mixes:

\begin{verbatim}
project intention
current architecture
configuration
operational behavior
test evidence
historical snapshot
\end{verbatim}

DDTA must classify those meanings before documenting them.

\subsection{4. A1 - From technology stack to project
problem}\label{4-a1---from-technology-stack-to-project-problem}

OR2 starts with a concrete stack: EfficientNet-B4, Qwen2-VL, ChromaDB,
BioMistral, B4 integration and a four-stage pipeline.

A literal reconstruction would incorrectly make those technologies the
project problem.

The stabilized framing instead became:

\begin{quote}
DermaTriage supports early triage of dermatological cases concerning
potentially oncological skin lesions, using available case information
to determine a priority of care and support routing toward appropriate
specialist assistance.
\end{quote}

\subsubsection{Important unresolved
boundary}\label{important-unresolved-boundary}

The source includes pathology/diagnostic outputs, but the documentation
does not sufficiently establish that DermaTriage owns definitive
clinical diagnosis authority. That meaning remains explicitly unresolved
rather than being hidden by generic wording.

\subsection{5. A2 - MacroRequirement
discovery}\label{5-a2---macrorequirement-discovery}

\subsubsection{5.1 MR-01 versus MR-02}\label{51-mr-01-versus-mr-02}

Source output bundles urgency/P-scale, specialist and SLA information.
Co-location does not prove one responsibility.

Split test:

\begin{verbatim}
Can priority/urgency change independently from specialist destination?
YES.
\end{verbatim}

Therefore:

\begin{itemize}
\tightlist
\item
  MR-01 owns triage urgency/priority;
\item
  MR-02 owns specialist-oriented routing indication.
\end{itemize}

MR-02 later STOPs because the source does not define enough downstream
routing semantics.

\subsubsection{5.2 MR-03 versus MR-04}\label{52-mr-03-versus-mr-04}

The source also bundles doctor review, prompt evolution and classifier
retraining into a ``self-learning'' story.

Split test:

\begin{verbatim}
Can clinical validation exist without model adaptation? YES.
Can adaptation policy change while review capture remains stable? YES.
\end{verbatim}

Therefore:

\begin{itemize}
\tightlist
\item
  MR-03 owns clinical-review result management;
\item
  MR-04 owns controlled adaptation using accumulated review evidence.
\end{itemize}

\texttt{MR-04\ dependsOn\ MR-03} because adaptation consumes the
clinical-review meaning.

\subsubsection{Methodology lesson}\label{methodology-lesson-1}

A lifecycle narrative is not automatically one MacroRequirement. Split
by autonomous responsibility, not by source heading.

\subsection{6. A3/A4 - MR-01: image fallback, P-scale and
SLA}\label{6-a3a4---mr-01-image-fallback-p-scale-and-sla}

\subsubsection{6.1 DEC-01 - no-image
continuity}\label{61-dec-01---no-image-continuity}

Source evidence says that when no image is available, symptom-only
scoring derives urgency from available B4 chatbot fields.

The material commitment surviving implementation substitution is:

\begin{verbatim}
image absence is not an absolute precondition for triage
\end{verbatim}

This supports DEC-01 and then FR-01.

\subsubsection{6.2 DEC-02 versus FR-02}\label{62-dec-02-versus-fr-02}

Source table:

\begin{verbatim}
HIGH + confidence > 0.85 -> P1
HIGH                     -> P2
MEDIUM                   -> P3
LOW                      -> P4
\end{verbatim}

The Decision is not the entire table. The material project commitment is
using P1-P4 as operational priority. The concrete conditional mapping is
an operational obligation and therefore becomes FR-02.

This separation prevents a Decision from becoming a disguised algorithm
table.

\subsubsection{6.3 SLA - precision without
sufficiency}\label{63-sla---precision-without-sufficiency}

The same source table lists 24h/48h/72h/7 days.

The review asked:

\begin{itemize}
\tightlist
\item
  what exactly must happen within the interval?
\item
  when does the timer start?
\item
  who owns compliance?
\item
  is the number a maximum, target or recommendation?
\end{itemize}

The source does not establish those semantics sufficiently.

Disposition:

\begin{verbatim}
GAP-DERMA-SLA-01
NOT SPECIFIED
\end{verbatim}

\subsubsection{Methodology lesson}\label{methodology-lesson-2}

A number can be precise while its meaning is underspecified.

\subsection{7. MR-03 - clinical review without stealing clinical
authority}\label{7-mr-03---clinical-review-without-stealing-clinical-authority}

A naive formulation such as:

\begin{verbatim}
DermaTriage validates its diagnoses clinically
\end{verbatim}

would incorrectly assign clinical judgment to the system.

The responsibility boundary established:

\begin{verbatim}
healthcare professional -> owns clinical judgment
DermaTriage             -> records/correlates review result
\end{verbatim}

DEC-03 therefore governs separation between original AI outcome and
clinical review, while FR-03 governs recording/correlation and
distinction between confirmation and correction.

The source field \texttt{agrees} was not promoted into FR-03 because a
boolean implementation representation is not automatically the semantic
definition of clinical review.

\subsection{8. MR-04 - first major holdout
pressure}\label{8-mr-04---first-major-holdout-pressure}

The architecture describes:

\begin{itemize}
\tightlist
\item
  prompt evolution every 10 corrections;
\item
  sliding window of 20 recent examples;
\item
  classifier retraining after 50 disagreement corrections;
\item
  \texttt{agrees\ ==\ False};
\item
  P-scale to HIGH/MEDIUM/LOW label mapping;
\item
  comparative acceptance criteria;
\item
  rollback beyond 5\% degradation;
\item
  fine-tuning details.
\end{itemize}

Treating all of this as ``continuous learning requirement'' would create
a monolithic and implementation-contaminated artifact.

The Decision family was therefore discovered before writing FRs.

\subsection{9. DEC-04 - numeric neutralization and accumulation
policy}\label{9-dec-04---numeric-neutralization-and-accumulation-policy}

Initial source phrasing:

\begin{verbatim}
Every 10 corrections -> prompt evolution
50 disagreement corrections -> classifier retraining
\end{verbatim}

Neutralization:

\begin{verbatim}
10 -> N
50 -> N
\end{verbatim}

What remains?

\begin{verbatim}
adaptation is not immediate per single review;
each path is activated after a governed accumulation condition is satisfied
\end{verbatim}

That surviving strategy is DEC-04.

\subsubsection{Why hardcoding 10/50 in FRs was
rejected}\label{why-hardcoding-1050-in-frs-was-rejected}

An early FR wording used exact literals. Review showed that the source
establishes current values but not their immutable normative status or
change authority.

Correct pattern:

\begin{verbatim}
FR semantic rule
    -> semantic threshold parameter N
    -> current documented binding 10 or 50
    -> realization configuration
\end{verbatim}

This became the basis for the Parameter Governance Boundary.

\subsection{10. DEC-09 / FR-06 - policy versus window-size
binding}\label{10-dec-09--fr-06---policy-versus-window-size-binding}

Source evidence says prompt evolution uses a sliding window of 20 most
recent examples.

Neutralizing \texttt{20\ -\textgreater{}\ N} leaves a material policy:

\begin{verbatim}
use bounded recent evidence rather than unrestricted history
\end{verbatim}

Therefore:

\begin{itemize}
\tightlist
\item
  DEC-09 owns the bounded-recent-evidence strategy;
\item
  FR-06 operationalizes selection of a bounded sliding window;
\item
  \texttt{PromptEvidenceWindowSize=N} is a semantic parameter;
\item
  current binding \texttt{N=20} is preserved without claiming immutable
  governance.
\end{itemize}

\subsubsection{Gap exposed}\label{gap-exposed}

The source does not fully define ordering, membership, underfill,
overlap/reuse, deduplication or scope. These stay in
\texttt{GAP-DERMA-PROMPT-WINDOW-01}.

\subsection{11. DEC-10 / FR-07 - semantic concept versus data-state
encoding}\label{11-dec-10--fr-07---semantic-concept-versus-data-state-encoding}

Source says classifier retraining uses corrections where:

\begin{verbatim}
agrees == False
\end{verbatim}

The semantic question is not ``must the field be named
\texttt{agrees}?''. The stable meaning is:

\begin{verbatim}
clinical disagreement with the relevant AI outcome qualifies as classifier-adaptation evidence
\end{verbatim}

Therefore:

\begin{verbatim}
ClinicianDisagreement
    -> governed semantic concept

agrees == False
    -> source-observed current data/state encoding
\end{verbatim}

This was the first clear example of the bridge:

\begin{verbatim}
semantic concept -> data/state encoding
\end{verbatim}

\subsubsection{Cross-MR gap}\label{cross-mr-gap}

MR-03 owns review-result semantics, while MR-04 owns its use as
adaptation evidence. The exact relation among correction, disagreement
and \texttt{agrees} is not fully defined, producing
\texttt{GAP-DERMA-REVIEW-DISAGREEMENT-BINDING-01}.

\subsection{12. DEC-11 / FR-08 - transformation is not a
parameter}\label{12-dec-11--fr-08---transformation-is-not-a-parameter}

Source mapping:

\begin{verbatim}
P1 -> HIGH
P2 -> HIGH
P3 -> MEDIUM
P4 -> LOW
\end{verbatim}

Replacing \texttt{HIGH/MEDIUM/LOW} by generic variables would destroy
the governed meaning. The mapping is therefore the functional
transformation itself, not a configuration parameter.

This produced FR-08 as one coherent requirement with four normative
clauses.

\subsubsection{Cross-MR ownership}\label{cross-mr-ownership}

The P-scale domain belongs to MR-01. MR-04 consumes the corrected
P-scale as classifier-supervision input. Consumption does not create a
second parent and does not automatically establish
\texttt{MR-04\ dependsOn\ MR-01}.

\subsection{13. DEC-06 / FR-09 - qualification is not
deployment}\label{13-dec-06--fr-09---qualification-is-not-deployment}

Source evidence says a retrained candidate must satisfy comparative
quality conditions before replacement.

FR-09 therefore produces a semantic result:

\begin{verbatim}
qualified for adoption
\end{verbatim}

not:

\begin{verbatim}
automatically deployed
\end{verbatim}

The source does not sufficiently govern final promotion authority.
\texttt{GAP-DERMA-DEPLOY-01} remains open.

This distinction prevents an evaluation gate from silently becoming
production-deployment authority.

\subsection{14. DEC-07 - first non-security specialization
pressure}\label{14-dec-07---first-non-security-specialization-pressure}

The source uses three asymmetric quality conditions:

\begin{verbatim}
sensitivity must not worsen
false-low must not worsen
accuracy may degrade only within a bounded tolerance
\end{verbatim}

These do not define a new autonomous capability. They strengthen the
acceptable satisfaction of FR-09.

A5 disposition:

\begin{verbatim}
NO NEW FR
ROUTE TO A7
\end{verbatim}

The likely specialization meanings are:

\begin{itemize}
\tightlist
\item
  \texttt{SensitivityNonDegradation};
\item
  \texttt{FalseLowNonDegradation};
\item
  \texttt{LimitedAccuracyDegradation}.
\end{itemize}

\subsubsection{Thesis-boundary lesson}\label{thesis-boundary-lesson}

This holdout confirms that
\texttt{SpecializedRequirement\ {[}abstract{]}} was a sensible design:
non-security specializations can exist. But the thesis is about security
analysis over governed documentation. Therefore the quality meanings are
preserved without opening a full \texttt{QualityRequirement} subtype
research track.

\subsection{15. The two 5\% values - same literal, different
semantics}\label{15-the-two-5-values---same-literal-different-semantics}

The source reports \texttt{5\%} in two places:

\begin{itemize}
\tightlist
\item
  pre-adoption bounded accuracy degradation;
\item
  post-adoption rollback trigger.
\end{itemize}

The review deliberately did not merge them.

\begin{verbatim}
AccuracyDegradationTolerance = T
RollbackAccuracyDegradationThreshold = R

current T = 5%
current R = 5%

T == R ? NOT ESTABLISHED
\end{verbatim}

Lifecycle position and purpose define parameter identity, not literal
equality.

The source also does not sufficiently establish whether \texttt{5\%}
means relative percent or percentage points in the governed contract,
even if current test documents perform a concrete calculation.

\subsection{16. DEC-08 / FR-10 - reversibility without inventing
monitoring}\label{16-dec-08--fr-10---reversibility-without-inventing-monitoring}

R4 stabilized DEC-08 as post-adoption reversibility.

The source currently describes ``automatic rollback'', but that wording
appears inside current realization/training documentation. The governed
requirement extracted is therefore narrower:

\begin{verbatim}
an already adopted materially degrading classifier adaptation must be revocable by restoring a previous acceptable state/version
\end{verbatim}

The review did \textbf{not} infer:

\begin{itemize}
\tightlist
\item
  continuous monitoring;
\item
  hourly evaluation;
\item
  automatic rollback as a permanent policy;
\item
  admin approval;
\item
  exact target-version algorithm.
\end{itemize}

Those meanings remain in \texttt{GAP-DERMA-ROLLBACK-BINDING-01} unless
later authority establishes them.

\subsection{17. DEC-05 cumulative review - why FR-11 was
necessary}\label{17-dec-05-cumulative-review---why-fr-11-was-necessary}

At first it appeared that DEC-05 might STOP without its own FR because
the prompt and classifier branches already had separate FRs.

Counterexample:

\begin{verbatim}
prompt threshold satisfied
-> system activates prompt evolution AND classifier adaptation

classifier threshold not satisfied
\end{verbatim}

FR-04 can still pass. FR-05 can still pass whenever its own condition
eventually occurs. Yet DEC-05 is violated because the paths are
cross-triggered.

This demonstrates:

\begin{verbatim}
all current FRs pass
while Decision still fails
\end{verbatim}

Therefore FR-11 was required to govern lifecycle independence.

\subsubsection{New high-confidence guide
test}\label{new-high-confidence-guide-test}

Before STOP below a Decision ask:

\begin{verbatim}
Can all current downstream Requirements be satisfied while this Decision is still violated?
\end{verbatim}

If yes, downstream coverage is incomplete.

\subsection{18. A5 cumulative closure}\label{18-a5-cumulative-closure}

Final FR set through A5:

\begin{longtable}[]{@{}lll@{}}
\toprule\noalign{}
FR & Parent & Meaning \\
\midrule\noalign{}
\endhead
\bottomrule\noalign{}
\endlastfoot
FR-01 & DEC-01 & symptom-based urgency when image unavailable \\
FR-02 & DEC-02 & urgency/confidence -\textgreater{} P-scale mapping \\
FR-03 & DEC-03 & correlated clinical review distinct from original
outcome \\
FR-04 & DEC-04 & prompt-path activation at semantic accumulation
threshold \\
FR-05 & DEC-04 & classifier-path activation at semantic
qualifying-evidence threshold \\
FR-06 & DEC-09 & bounded recent evidence for prompt evolution \\
FR-07 & DEC-10 & clinical disagreement qualification \\
FR-08 & DEC-11 & corrected P-scale -\textgreater{} classifier
supervision target \\
FR-09 & DEC-06 & comparative pre-adoption qualification \\
FR-10 & DEC-08 & post-adoption rollback capability \\
FR-11 & DEC-05 & lifecycle independence between adaptation paths \\
\end{longtable}

DEC-07 routes to A7 rather than producing a fake FR.

A5 cumulative disposition:

\begin{verbatim}
PASS / CLOSED
A4 REOPEN: NO
\end{verbatim}

\subsection{19. Methodology findings accumulated through
A5}\label{19-methodology-findings-accumulated-through-a5}

\subsubsection{GI-13 - Numeric neutralization before parameter
classification}\label{gi-13---numeric-neutralization-before-parameter-classification}

A concrete value should not determine document type merely because it is
numeric.

\subsubsection{GI-14 - Cross-MR/reference is not
ownership}\label{gi-14---cross-mrreference-is-not-ownership}

A requirement may consume governed meaning from another owner without a
second parent.

\subsubsection{GI-21 - Information/realization
bridge}\label{gi-21---informationrealization-bridge}

Observed bridge taxonomy:

\begin{verbatim}
semantic concept -> data/state encoding
semantic parameter -> configuration value
semantic reference -> interface identifier
semantic transformation -> realization encoding
\end{verbatim}

\subsubsection{GI-22 - Requirement identity versus number of
clauses/test
rows}\label{gi-22---requirement-identity-versus-number-of-clausestest-rows}

One coherent obligation may contain multiple normative clauses and
multiple test cases.

\subsubsection{GI-23 - Shared vocabulary does not prove cross-FR
binding}\label{gi-23---shared-vocabulary-does-not-prove-cross-fr-binding}

Lexical equality is not enough to establish producer/consumer identity.

\subsubsection{GI-24 - Semantic parameter versus concrete
binding}\label{gi-24---semantic-parameter-versus-concrete-binding}

Preserve symbolic semantic parameter separately from current concrete
value.

\subsubsection{GI-25 - Parameter Governance
Boundary}\label{gi-25---parameter-governance-boundary}

Semantic-owner and lifecycle classification must occur before deciding
whether a numeric value belongs in the governed semantic layer.

\subsubsection{A5 counterexample test - candidate new guide
rule}\label{a5-counterexample-test---candidate-new-guide-rule}

Before Decision STOP, verify that the Decision cannot still be violated
while all current descendants pass.

\subsection{20. What the source documentation did
well}\label{20-what-the-source-documentation-did-well}

DermaTriage source documentation contains unusually useful concrete
evidence:

\begin{itemize}
\tightlist
\item
  explicit four-stage pipeline and adaptation layer;
\item
  endpoint inventory;
\item
  explicit fallback behavior;
\item
  clinical-review endpoints;
\item
  threshold/window values;
\item
  label mapping;
\item
  comparative quality criteria;
\item
  rollback evidence;
\item
  privacy/data-handling statements;
\item
  test and training evidence.
\end{itemize}

This concreteness made semantic classification possible.

\subsection{21. What the source documentation leaves
ambiguous}\label{21-what-the-source-documentation-leaves-ambiguous}

The same documentation mixes stable project meaning with current
implementation snapshot. Important examples:

\begin{itemize}
\tightlist
\item
  whether current numeric values are governed constants or
  configuration;
\item
  whether deployment is automatic after qualification;
\item
  exact acceptance and rollback quantitative semantics;
\item
  complete clinical-review state model;
\item
  specialist-routing semantics;
\item
  SLA meaning;
\item
  diagnostic authority;
\item
  exact data/interface bindings and missing-value behavior.
\end{itemize}

DDTA adds value by making these boundaries explicit instead of hiding
them inside prose/tables.

\subsection{22. Security-related source evidence intentionally
deferred}\label{22-security-related-source-evidence-intentionally-deferred}

The source also contains:

\begin{itemize}
\tightlist
\item
  X-API-Key protection on many administrative endpoints;
\item
  B4 JWT authentication;
\item
  anonymized patient IDs in research data;
\item
  no raw images in the RAG CSV;
\item
  B4 patient uploads processed in memory and not permanently stored in
  the research dataset.
\end{itemize}

These facts are preserved for A8. They are not yet promoted in this
study to final SecurityRequirements because the security-property scope
and governance must be reviewed explicitly.

This preserves the thesis order:

\begin{verbatim}
document project meaning first
-> review existing governed security properties
-> close documentation
-> build BA
-> run STRIDE/STRIDE-AI
-> feed accepted security findings back through governance
\end{verbatim}

\subsection{23. Why BA is still
blocked}\label{23-why-ba-is-still-blocked}

The accepted DermaTriage BA must not begin until A12 promotion readiness
passes. A limited pre-A7 consumer probe may later be used to test
whether the BA can preserve generic specialization information, but it
must not decide project meaning or document types.

\subsection{24. Next research steps}\label{24-next-research-steps}

\begin{verbatim}
A6 project-wide Requirement split
-> PRE-A7 scope + Documentation<->BA consistency check
-> A7 specialization review
   -> preserve non-security specializations outside thesis subtype scope
-> A8 existing governed SecurityRequirement review
-> A9 cross-MR
-> A10 terminology/bindings
-> A11 regression
-> A12 promotion
-> accepted DermaTriage BA
-> BA stabilization/regression
-> STRIDE / STRIDE-AI
-> governed security feedback cycle
\end{verbatim}

\subsection{25. Core lesson of the holdout so
far}\label{25-core-lesson-of-the-holdout-so-far}

The most useful result is not that DDTA can restate the architecture
document. It is that the method repeatedly separates four things that
ordinary documentation tends to mix:

\begin{verbatim}
what must remain true for project meaning to remain the same
what is a current concrete binding
what is current realization evidence
what is genuinely not specified
\end{verbatim}

That separation is what makes the later Base Analysis and security
analysis traceable without allowing the analysis method to rewrite
project truth.

\end{document}
